\documentclass[11pt]{article}

\usepackage[margin=1in]{geometry}
\usepackage{amsmath}
\usepackage{amssymb}
\usepackage{booktabs}
\usepackage{longtable}

\newcommand{\dd}{\mathrm{d}}
\newcommand{\diag}{\operatorname{diag}}
\newcommand{\sgn}{\operatorname{sgn}}
\newcommand{\sat}{\operatorname{sat}}

\title{Minilink Dynamics Catalog Equations of Motion}
\author{Minilink}
\date{\today}

\begin{document}
\maketitle

\section*{Scope and conventions}
This document records the equations of motion for the clean-slate
\texttt{minilink.dynamics.catalog} systems migrated from Pyro.  The equations
were checked against the corresponding Pyro catalog classes by numerical
comparison of state derivatives, state-space matrices, and mechanical model
terms.

Mechanical systems use
\begin{equation}
    H(q)\ddot q + C(q,\dot q)\dot q + d(q,\dot q,u,t) + g(q)
    = \tau(q,\dot q,u,t).
\end{equation}
Generalized mechanical systems use
\begin{equation}
    M(q)\dot v + C(q,v)v + d(q,v,u,t) + g(q) = \tau(q,v,u,t),
    \qquad
    \dot q = N(q)v.
\end{equation}
Linear state-space systems use
\begin{equation}
    \dot x = A x + B u,
    \qquad
    y = Cx + Du.
\end{equation}

\section{Equation systems}

\subsection{SimpleIntegrator}
\begin{equation}
    x = p, \qquad u = v, \qquad \dot p = v.
\end{equation}

\subsection{DoubleIntegrator}
\begin{equation}
    x = \begin{bmatrix} p \\ v \end{bmatrix},
    \qquad
    \dot x = \begin{bmatrix} v \\ u \end{bmatrix},
    \qquad
    y = p.
\end{equation}

\subsection{TripleIntegrator}
\begin{equation}
    x = \begin{bmatrix} p \\ v \\ f \end{bmatrix},
    \qquad
    \dot x = \begin{bmatrix} v \\ f \\ u \end{bmatrix},
    \qquad
    y = p.
\end{equation}

\subsection{VanderPol}
\begin{equation}
    x = \begin{bmatrix} y \\ \dot y \end{bmatrix},
    \qquad
    \dot x =
    \begin{bmatrix}
        \dot y \\
        -y + \mu \dot y(1-y^2)
    \end{bmatrix}.
\end{equation}

\subsection{TransferFunction}
For a SISO transfer function
\begin{equation}
    G(s) = \frac{b_0s^r + \cdots + b_r}{a_0s^n + \cdots + a_n},
\end{equation}
the implementation uses SciPy's controllable state-space realization:
\begin{equation}
    (A,B,C,D) = \operatorname{tf2ss}(b,a),
    \qquad
    \dot x = Ax + Bu,
    \qquad
    y = Cx + Du.
\end{equation}

\section{Mass-spring-damper}

\subsection{SingleMass}
\begin{equation}
    x = \begin{bmatrix} p \\ v \end{bmatrix},
    \qquad
    A =
    \begin{bmatrix}
        0 & 1 \\
        -k/m & -b/m
    \end{bmatrix},
    \qquad
    B = \begin{bmatrix}0 \\ 1/m\end{bmatrix}.
\end{equation}

\subsection{FloatingSingleMass}
\begin{equation}
    \dot p = v,
    \qquad
    \dot v = \frac{u - b v}{m}.
\end{equation}
This is \texttt{SingleMass} with \(k=0\).

\subsection{TwoMass}
\begin{align}
    \dot p_1 &= v_1, &
    \dot p_2 &= v_2, \\
    \dot v_1 &=
        \frac{-(k_1+k_2)p_1 + k_2p_2 - b_1v_1}{m_1}, &
    \dot v_2 &=
        \frac{k_2p_1 - k_2p_2 - b_2v_2 + u}{m_2}.
\end{align}
The output is \(p_2\) by default, or \(p_1\) when requested.

\subsection{FloatingTwoMass}
\begin{equation}
    \text{\texttt{FloatingTwoMass}: same as \texttt{TwoMass} with } k_1=0.
\end{equation}

\subsection{ThreeMass}
\begin{align}
    \dot p_1 &= v_1, &
    \dot p_2 &= v_2, &
    \dot p_3 &= v_3, \\
    \dot v_1 &=
        \frac{-(k_1+k_2)p_1 + k_2p_2 - b_1v_1}{m_1}, \\
    \dot v_2 &=
        \frac{k_2p_1 -(k_2+k_3)p_2 + k_3p_3 - b_2v_2}{m_2}, \\
    \dot v_3 &=
        \frac{k_3p_2-k_3p_3-b_3v_3+u}{m_3}.
\end{align}
The default Pyro output is \(p_2\); \(p_1,p_2,p_3\) are selectable.

\subsection{FloatingThreeMass}
\begin{equation}
    \text{\texttt{FloatingThreeMass}: same as \texttt{ThreeMass} with } k_1=0,
    \text{ and default output } p_3.
\end{equation}

\section{Pendulum and cart-pole}

\subsection{Pendulum}
With \(q=\theta\),
\begin{equation}
    H = m_1l_{c1}^2 + I_1,
    \qquad
    C = 0,
    \qquad
    g = m_1g_0l_{c1}\sin\theta,
    \qquad
    d = d_1\dot\theta,
    \qquad
    \tau = u.
\end{equation}

\subsection{InvertedPendulum}
\begin{equation}
    H\ddot\theta + d_1\dot\theta - m_1g_0l_{c1}\sin\theta = u.
\end{equation}

\subsection{TwoIndependentPendulums}
For \(q=[\theta_1,\theta_2]^T\),
\begin{equation}
    H = hI_2,\qquad
    h=m_1l_{c1}^2+I_1,\qquad
    C=0,\qquad
    g=m_1g_0l_{c1}\sin(q),\qquad
    d=d_1\dot q,\qquad
    \tau=u.
\end{equation}

\subsection{DoublePendulum}
Let \(q=[q_1,q_2]^T\), \(c_2=\cos q_2\), \(s_2=\sin q_2\), and
\(s_{12}=\sin(q_1+q_2)\).  The Pyro double-pendulum convention is
\begin{align}
    H_{11} &=
        m_1l_{c1}^2 + I_1
        + m_2(l_1^2+l_{c2}^2+2l_1l_{c2}c_2) + I_2, \\
    H_{12} &= H_{21} =
        m_2l_{c2}^2 + m_2l_1l_{c2}c_2 + I_2, \\
    H_{22} &= m_2l_{c2}^2 + I_2,
\end{align}
\begin{equation}
    C =
    \begin{bmatrix}
        -h\dot q_2 & -h(\dot q_1+\dot q_2) \\
        h\dot q_1 & 0
    \end{bmatrix},
    \qquad
    h=m_2l_1l_{c2}s_2,
\end{equation}
\begin{equation}
    g =
    \begin{bmatrix}
        -[(m_1l_{c1}+m_2l_1)g_0]\sin q_1 - m_2l_{c2}g_0s_{12} \\
        -m_2l_{c2}g_0s_{12}
    \end{bmatrix},
    \qquad
    d=\diag(d_1,d_2)\dot q,
    \qquad
    \tau=u.
\end{equation}

\subsection{Acrobot}
\begin{equation}
    H,C,g,d \text{ are the same as \texttt{DoublePendulum}}, \qquad
    \tau = \begin{bmatrix}0\\u\end{bmatrix}.
\end{equation}

\subsection{RotatingCartPole}
For \(q=[q_1,q_2]^T\),
\begin{equation}
    H =
    \begin{bmatrix}
        m_2l_1^2+I_1 & m_2l_1l_2\cos q_2 \\
        m_2l_1l_2\cos q_2 & m_2l_2^2+I_2
    \end{bmatrix},
\end{equation}
\begin{equation}
    C =
    \begin{bmatrix}
        0 & -m_2l_1l_2\sin q_2\dot q_2 \\
        0 & 0
    \end{bmatrix},
    \qquad
    g =
    \begin{bmatrix}
        0 \\
        -m_2g_0l_2\sin q_2
    \end{bmatrix},
    \qquad
    d=\diag(d_1,d_2)\dot q,
    \qquad
    \tau=u.
\end{equation}

\subsection{UnderactuatedRotatingCartPole}
\begin{equation}
    H,C,g,d \text{ are the same as \texttt{RotatingCartPole}}, \qquad
    \tau=\begin{bmatrix}u\\0\end{bmatrix}.
\end{equation}

\subsection{CartPole and JaxCartPole}
For \(q=[x,\theta]^T\), \(\dot q=[\dot x,\dot\theta]^T\),
\begin{equation}
    H =
    \begin{bmatrix}
        m_1+m_2 & m_2l_{cg}\cos\theta \\
        m_2l_{cg}\cos\theta & m_2l_{cg}^2
    \end{bmatrix},
    \qquad
    C =
    \begin{bmatrix}
        0 & -m_2l_{cg}\sin\theta\,\dot\theta \\
        0 & 0
    \end{bmatrix},
\end{equation}
\begin{equation}
    g =
    \begin{bmatrix}
        0\\m_2g_0l_{cg}\sin\theta
    \end{bmatrix},
    \qquad
    d=0,
    \qquad
    \tau=\begin{bmatrix}u\\0\end{bmatrix}.
\end{equation}
\texttt{JaxCartPole} has the same equations and JAX-traceable array builders.

\section{Vehicles}

\subsection{KinematicBicycle, KinematicCar, and UdeSRacecar}
\begin{equation}
    \dot x = v\cos\theta,
    \qquad
    \dot y = v\sin\theta,
    \qquad
    \dot\theta = \frac{v\tan\delta}{L}.
\end{equation}
\texttt{KinematicCar} and \texttt{UdeSRacecar} use the same model with their
own wheelbase \(L=a+b\).

\subsection{ConstantSpeedKinematicCar}
\begin{equation}
    \dot x = v_0\cos\theta,
    \qquad
    \dot y = v_0\sin\theta,
    \qquad
    \dot\theta = \frac{v_0\tan\delta}{L}.
\end{equation}

\subsection{HolonomicMobileRobot}
\begin{equation}
    x=\begin{bmatrix}p_x\\p_y\end{bmatrix},
    \qquad
    u=\begin{bmatrix}v_x\\v_y\end{bmatrix},
    \qquad
    \dot x=u.
\end{equation}

\subsection{HolonomicMobileRobot3D}
\begin{equation}
    x=\begin{bmatrix}p_x\\p_y\\p_z\end{bmatrix},
    \qquad
    u=\begin{bmatrix}v_x\\v_y\\v_z\end{bmatrix},
    \qquad
    \dot x=u.
\end{equation}

\subsection{LongitudinalFrontWheelDriveCarWithWheelSlipInput}
Let \(x=[p,v]^T\), \(s\) be wheel slip, and
\begin{equation}
    \mu(s) = \mu_{\max}\left(\frac{2}{1+\exp(-\mu_{\mathrm{slope}}s)}-1\right).
\end{equation}
With
\begin{equation}
    r_y=\frac{y_c}{L},\qquad r_r=\frac{x_c}{L},\qquad
    D(v)=\frac{1}{2}\rho C_dA\,v|v|,
\end{equation}
the model is
\begin{equation}
    \dot p=v,
    \qquad
    \dot v =
    \frac{\mu(s)mg_0r_r-D(v)}{m(1+\mu(s)r_y)}.
\end{equation}

\subsection{LongitudinalFrontWheelDriveCarWithTorqueInput}
For \(x=[p,v,\omega,\phi]^T\) and wheel torque \(T\),
\begin{equation}
    s =
    \sat_{[-0.5,0.5]}\left(\frac{R\omega-v}{|v|+\epsilon}\right),
    \qquad
    a = \frac{\mu(s)mg_0r_r-D(v)}{m(1+\mu(s)r_y)}.
\end{equation}
\begin{equation}
    \dot p=v,
    \qquad
    \dot v=a,
    \qquad
    \dot\omega=\frac{T-R(ma+D(v))}{I_w},
    \qquad
    \dot\phi=\omega.
\end{equation}

\subsection{QuarterCarOnRoughTerrain}
With state \(x_s=[\dot y,y,x]^T\), prescribed forward speed \(v_x\), and
\begin{equation}
    z(x)=\sum_i a_i\sin(w_i(x-\phi_i)),\qquad
    z'(x)=\sum_i a_iw_i\cos(w_i(x-\phi_i)),
\end{equation}
the dynamics are
\begin{equation}
    \ddot y =
    \frac{u-k(y-z(x))-b(\dot y-z'(x))}{m},
    \qquad
    \dot x=v_x.
\end{equation}

\subsection{MountainCar}
The car is constrained to \(z(q)=a\cos(wq)\):
\begin{align}
    H(q)&=m(1+z'(q)^2),
    &
    C(q,\dot q)&=mz'(q)z''(q)\dot q,\\
    B(q)&=\sqrt{1+z'(q)^2},
    &
    g(q)&=mg_0z'(q),
    &
    d&=0.
\end{align}
Thus \(H\ddot q+C\dot q+g=B(q)u\).

\subsection{DynamicBicycle and tire laws}
For \(q=[x,y,\theta]^T\), \(v=[v_x,v_y,r]^T\),
\begin{equation}
    \dot q =
    \begin{bmatrix}
        \cos\theta & -\sin\theta & 0\\
        \sin\theta & \cos\theta & 0\\
        0&0&1
    \end{bmatrix}v,
    \qquad
    M=\diag(m,m,I_z),
\end{equation}
\begin{equation}
    C =
    \begin{bmatrix}
        0&-mr&0\\
        mr&0&0\\
        0&0&0
    \end{bmatrix}.
\end{equation}
For steering input \(\delta\) and rear wheel speed \(\omega_r\),
\begin{align}
    v_{xf} &= \cos\delta\,v_x+\sin\delta\,(v_y+ar),\\
    v_{yf} &= -\sin\delta\,v_x+\cos\delta\,(v_y+ar),\\
    v_{xr} &= v_x,\qquad v_{yr}=v_y-br,\qquad
    \omega_f = v_{xf}/R_f.
\end{align}
The tire slip variables are
\begin{equation}
    \alpha=-\arctan\left(\frac{v_y}{|v_x|+\epsilon}\right),
    \qquad
    \kappa=\frac{\omega R-v_x}{|v_x|+\epsilon}.
\end{equation}
The linear tire law is
\begin{equation}
    F_x=C_\kappa\kappa,\qquad F_y=C_\alpha\alpha,
    \qquad
    \|(F_x,F_y)\|\leq \mu F_z,
\end{equation}
with friction-circle saturation.  The Pacejka law is
\begin{equation}
    F = D F_z \sin\left(C\arctan(Bs-E(Bs-\arctan(Bs)))\right)
\end{equation}
applied to \(s=\kappa\) and \(s=\alpha\).  Body-frame front forces are
\begin{equation}
    F_{xf}^{b}=F_{xf}\cos\delta-F_{yf}\sin\delta,
    \qquad
    F_{yf}^{b}=F_{xf}\sin\delta+F_{yf}\cos\delta.
\end{equation}
The left-side load is
\begin{equation}
    d=-
    \begin{bmatrix}
        F_{xf}^{b}+F_{xr}-\frac{1}{2}\rho C_dA v_x|v_x|\\
        F_{yf}^{b}+F_{yr}\\
        aF_{yf}^{b}-bF_{yr}
    \end{bmatrix},
    \qquad
    \tau=0.
\end{equation}

\section{Aerial systems}

\subsection{Drone2D}
For \(q=[x,y,\theta]^T\),
\begin{equation}
    H=\diag(m,m,I),
    \qquad
    C=0,
    \qquad
    g=\begin{bmatrix}0\\mg_0\\0\end{bmatrix},
\end{equation}
\begin{equation}
    d=
    \begin{bmatrix}
        c_d\dot x|\dot x|+0.01\dot x\\
        c_d\dot y|\dot y|+0.01\dot y\\
        0.01\dot\theta
    \end{bmatrix},
    \qquad
    B(q)=
    \begin{bmatrix}
        -\sin\theta & -\sin\theta\\
        \cos\theta & \cos\theta\\
        -l & l
    \end{bmatrix},
    \qquad
    \tau=B(q)\begin{bmatrix}T_1\\T_2\end{bmatrix}.
\end{equation}

\subsection{Drone2DWithSideThruster}
\begin{equation}
    B_{\mathrm{side}}(q)=
    \begin{bmatrix}
        -\sin\theta & -\sin\theta & \cos\theta\\
        \cos\theta & \cos\theta & \sin\theta\\
        -l & l & 0
    \end{bmatrix},
    \qquad
    \tau=B_{\mathrm{side}}(q)
    \begin{bmatrix}T_1\\T_2\\T_L\end{bmatrix}.
\end{equation}

\subsection{SpeedControlledDrone2D}
\begin{equation}
    x=\begin{bmatrix}p_x\\p_y\end{bmatrix},
    \qquad
    \dot x=\begin{bmatrix}u_x\\u_y\end{bmatrix}.
\end{equation}

\subsection{ConstantSpeedHelicopterTunnel}
With state \(x_s=[\dot y,y,x]^T\),
\begin{equation}
    \ddot y = \frac{u}{m},
    \qquad
    \dot x=v_x.
\end{equation}

\subsection{Rocket}
\begin{equation}
    H=\diag(m,m,I),
    \qquad
    g=\begin{bmatrix}0\\mg_0\\0\end{bmatrix},
    \qquad
    d=
    \begin{bmatrix}
        c_d\dot x|\dot x|+0.01\dot x\\
        c_d\dot y|\dot y|+0.01\dot y\\
        0.01\dot\theta
    \end{bmatrix},
\end{equation}
\begin{equation}
    \tau =
    T
    \begin{bmatrix}
        -\sin(\theta+\delta)\\
        \cos(\theta+\delta)\\
        -y_{cg}\sin\delta
    \end{bmatrix}.
\end{equation}

\subsection{Plane2D}
For \(q=[x,y,\theta]^T\), \(\dot q=[v_x,v_y,\omega]^T\),
\begin{equation}
    H=\diag(m,m,I),
    \qquad
    g=\begin{bmatrix}0\\mg_0\\0\end{bmatrix},
    \qquad
    \tau=T\begin{bmatrix}\cos\theta\\\sin\theta\\0\end{bmatrix}.
\end{equation}
Let
\begin{equation}
    V=\sqrt{v_x^2+v_y^2},
    \qquad
    \gamma=\arctan2(v_y,v_x),
    \qquad
    \alpha=\theta-\gamma.
\end{equation}
The lift, drag, and moment laws are
\begin{align}
    C_l(\alpha) &= \sin(2\alpha) + \begin{cases}4\alpha,&|\alpha|<\alpha_s\\0,&\text{otherwise,}\end{cases}\\
    C_d(\alpha) &= C_{d0} + 1-\cos(2\alpha)
      + \begin{cases}C_l(\alpha)^2/(\pi e AR),&|\alpha|<\alpha_s\\0,&\text{otherwise,}\end{cases}\\
    C_m(\alpha) &= C_{m0}.
\end{align}
With dynamic pressure \(\bar q=\frac{1}{2}\rho V^2\),
\begin{equation}
    L_w=\bar qS_wC_l(\alpha),\quad
    D_w=\bar qS_wC_d(\alpha),\quad
    M_w=\bar qS_wc_wC_m(\alpha),
\end{equation}
and the same expressions for the tail using \(S_t,c_t,\alpha+\delta\).  The
world-frame aerodynamic load is
\begin{equation}
    d =
    -R(\gamma)
    \begin{bmatrix}
        -(D_w+D_t)\\
        L_w+L_t\\
        M_w+M_t-l_w(L_w\cos\alpha+D_w\sin\alpha)
            -l_t(L_t\cos\alpha+D_t\sin\alpha)
    \end{bmatrix}.
\end{equation}

\section{Marine systems}

\subsection{Boat2D}
For \(q=[x,y,\theta]^T\), body velocity \(v=[u_b,v_b,r]^T\),
\begin{equation}
    M=\diag(m,m,I),
    \qquad
    N(q)=
    \begin{bmatrix}
        \cos\theta&-\sin\theta&0\\
        \sin\theta&\cos\theta&0\\
        0&0&1
    \end{bmatrix},
\end{equation}
\begin{equation}
    C=
    \begin{bmatrix}
        0&-mr&0\\
        mr&0&0\\
        0&0&0
    \end{bmatrix},
    \qquad
    B=
    \begin{bmatrix}
        1&0\\
        0&1\\
        0&-l_t
    \end{bmatrix},
    \qquad
    \tau=B\begin{bmatrix}T_x\\T_y\end{bmatrix}.
\end{equation}
For relative velocity \(v_r\),
\begin{equation}
    d(v_r)=D_{\ell}v_r+
    \begin{bmatrix}
        -\frac{1}{2}\rho A_f C_x(\alpha)\|v_{r,xy}\|^2\\
        -\frac{1}{2}\rho A_l C_y(\alpha)\|v_{r,xy}\|^2\\
        -\frac{1}{2}\rho A_l L C_m(\alpha)\|v_{r,xy}\|^2
        +N_{\max}\rho A_lL|r|r
    \end{bmatrix},
\end{equation}
where \(\alpha=-\arctan2(v_{r,y},v_{r,x})\),
\begin{equation}
    C_x=-C_{x,\max}\cos\alpha|\cos\alpha|,
    \quad
    C_y=C_{y,\max}\sin\alpha|\sin\alpha|,
    \quad
    C_m=C_{m,\max}\sin(2\alpha).
\end{equation}
For \texttt{Boat2D}, \(v_r=v\).

\subsection{Boat2DWithCurrent}
With world current \(v_c=[v_{cx},v_{cy},0]^T\),
\begin{equation}
    v_r = v - N(q)^Tv_c,
\end{equation}
and the rest of the \texttt{Boat2D} equations are unchanged.

\section{Manipulators}

\subsection{SpeedControlledManipulator}
\begin{equation}
    x=q,\qquad u=\dot q_{\mathrm{cmd}},\qquad \dot q=u.
\end{equation}
The forward differential kinematics are
\begin{equation}
    \dot r = J(q)\dot q.
\end{equation}

\subsection{OneLinkManipulator}
\begin{equation}
    H=m_1l_{c1}^2+I_1,
    \qquad
    C=0,
    \qquad
    g=-m_1g_0l_{c1}\sin q,
    \qquad
    d=d_1\dot q,
    \qquad
    \tau=u.
\end{equation}
The end-effector kinematics are
\begin{equation}
    r(q)=\begin{bmatrix}l_1\sin q\\l_1\cos q\end{bmatrix},
    \qquad
    J(q)=\begin{bmatrix}l_1\cos q\\-l_1\sin q\end{bmatrix}.
\end{equation}

\subsection{TwoLinkManipulator}
Let \(q=[q_1,q_2]^T\), \(c_2=\cos q_2\), \(s_2=\sin q_2\), and
\(c_{12}=\cos(q_1+q_2)\), \(s_{12}=\sin(q_1+q_2)\).  The inertia matrix is
\begin{align}
    H_{11} &= m_1l_{c1}^2+I_1
      +m_2(l_1^2+l_{c2}^2+2l_1l_{c2}c_2)+I_2,\\
    H_{12} &= H_{21}=m_2l_{c2}^2+m_2l_1l_{c2}c_2+I_2,\\
    H_{22} &= m_2l_{c2}^2+I_2.
\end{align}
\begin{equation}
    C=
    \begin{bmatrix}
        -h\dot q_2&-h(\dot q_1+\dot q_2)\\
        h\dot q_1&0
    \end{bmatrix},
    \quad
    h=m_2l_1l_{c2}s_2,
    \quad
    d=\diag(d_1,d_2)\dot q.
\end{equation}
\begin{equation}
    g=
    \begin{bmatrix}
        -(m_1l_{c1}+m_2l_1)g_0\sin q_1-m_2l_{c2}g_0s_{12}\\
        -m_2l_{c2}g_0s_{12}
    \end{bmatrix},
    \qquad
    \tau=u.
\end{equation}
\begin{equation}
    r(q)=
    \begin{bmatrix}
        l_1\sin q_1+l_2s_{12}\\
        l_1\cos q_1+l_2c_{12}
    \end{bmatrix}.
\end{equation}

\subsection{ThreeLinkManipulator3D}
Let \(q=[q_1,q_2,q_3]^T\), \(c_i=\cos q_i\), \(s_i=\sin q_i\), and
\(c_{23}=\cos(q_2+q_3)\), \(s_{23}=\sin(q_2+q_3)\).  With
\(\ell=\ell_2\), \(r_1=l_{c2}\), and \(r_2=l_{c3}\),
\begin{align}
    H_{11} &=
        I_{2y}s_2^2 + I_{3y}s_{23}^2 + I_{1z}
        + I_{2z}c_2^2 + I_{3z}c_{23}^2
        + m_2(r_1c_2)^2 + m_3(\ell c_2+r_2c_{23})^2,\\
    H_{22} &=
        I_{2x}+I_{3x}+m_3\ell^2+m_2r_1^2+m_3r_2^2
        +2m_3\ell r_2c_3,\\
    H_{23} &= H_{32}=I_{3x}+m_3r_2^2+m_3\ell r_2c_3,\\
    H_{33} &= I_{3x}+m_3r_2^2,
\end{align}
and all other entries of \(H\) are zero.  The Coriolis matrix uses
\begin{align}
    T_{112} &=
        (I_{2y}-I_{2z}-m_2r_1^2)c_2s_2
        +(I_{3y}-I_{3z})c_{23}s_{23}
        -m_3(\ell c_2+r_2c_{23})(\ell s_2+r_2s_{23}),\\
    T_{113} &=
        (I_{3y}-I_{3z})c_{23}s_{23}
        -m_3r_2s_{23}(\ell c_2+r_2c_{23}),\\
    T_{223} &=-\ell m_3r_2s_3,
    \qquad
    T_{322}=\ell m_3r_2s_3.
\end{align}
\begin{equation}
    C =
    \begin{bmatrix}
        T_{112}\dot q_2+T_{113}\dot q_3&T_{112}\dot q_1&T_{113}\dot q_1\\
        -T_{112}\dot q_1&T_{223}\dot q_3&T_{223}(\dot q_2+\dot q_3)\\
        -T_{113}\dot q_1&T_{322}\dot q_2&0
    \end{bmatrix}.
\end{equation}
\begin{equation}
    g =
    \begin{bmatrix}
        0\\
        -(m_2g_0l_{c2}+m_3g_0l_2)c_2-m_3g_0l_{c3}c_{23}\\
        -m_3g_0l_{c3}c_{23}
    \end{bmatrix},
    \qquad
    d=\diag(d_1,d_2,d_3)\dot q,
    \qquad
    \tau=u.
\end{equation}

\subsection{FiveLinkPlanarManipulator}
Pyro's five-link class defines kinematics and otherwise uses the default
fully actuated mechanical model:
\begin{equation}
    H=I_5,\qquad C=0,\qquad g=0,\qquad d=0,\qquad \tau=u.
\end{equation}
For link lengths \(l_i\) and absolute angles
\(\alpha_i=\sum_{j=1}^{i}q_j\),
\begin{equation}
    r(q)=
    \begin{bmatrix}
        \sum_{i=1}^{5} l_i\sin\alpha_i\\
        \sum_{i=1}^{5} l_i\cos\alpha_i
    \end{bmatrix},
\end{equation}
\begin{equation}
    J_{1j}=\sum_{i=j}^{5}l_i\cos\alpha_i,
    \qquad
    J_{2j}=-\sum_{i=j}^{5}l_i\sin\alpha_i.
\end{equation}

\end{document}
